% TOP_PROBLEM: {"problemId": "problem.provisional-top500-omission-waves-010-014.018", "canonicalTitle": "Hilbert's Tenth Problem over C(t) with named t", "releaseRank": 446, "primaryDomain": "number_theory_arithmetic_geometry", "primaryDomainLabel": "Number theory & arithmetic geometry", "displayStatus": "open", "releaseStatus": "open"}
% !TeX program = lualatex
% Definition and sources only; this file is not an LLM solution attempt.
\documentclass[11pt]{article}
\usepackage[margin=25mm]{geometry}
\usepackage{fontspec}
\setmainfont{Latin Modern Roman}
\usepackage{amsmath,amssymb,mathtools}
\usepackage{unicode-math}
\setmathfont{Latin Modern Math}
\usepackage{xurl,hyperref}
\hypersetup{colorlinks=true,urlcolor=blue}
\setcounter{secnumdepth}{0}
\setlength{\parindent}{0pt}
\setlength{\parskip}{5pt}
\setlength{\emergencystretch}{3em}
\begin{document}
\section*{\texorpdfstring{Hilbert's Tenth Problem over C(t) with named t}{Hilbert's Tenth Problem over C(t) with named t}}\label{446}
\subsection{Definitions and mathematical statement}
Let ${\mathbb{C}}(t)$ be the field of rational functions in an indeterminate $t$ with complex coefficients. The input is a finite polynomial
\[
 P(t,X_1,\ldots,X_n)\in{\mathbb{Z}}[t,X_1,\ldots,X_n],
\]
encoded by its integer coefficients and exponents. Decide whether there exist $x_1(t),\ldots,x_n(t)\in{\mathbb{C}}(t)$ with $P(t,x_1(t),\ldots,x_n(t))=0$ as a rational-function identity. The named element $t$ is part of the language, so this is not simply the algebraic-closedness problem for ${\mathbb{C}}$.

The target asks for one algorithm halting correctly for every such input, or an undecidability proof. The possible solutions may use arbitrary complex constants; those constants are not supplied as arbitrary uncomputable input data. The catalogue links this to the corresponding existential theory, with the language conventions of the cited source retained.
\par\smallskip\textit{Formulation and concept references:} \href{https://math.mit.edu/~poonen/papers/aws2003.pdf}{[S1]}.

\subsection{Short English statement}
Can an algorithm decide whether a polynomial equation with integer coefficients and a named variable t has solutions in the complex rational-function field?
\subsection{Sources}
\begin{itemize}
\item[S1]B. Poonen, Hilbert's tenth problem over rings of number-theoretic interest, Arizona Winter School notes, 2003.
\url{https://math.mit.edu/~poonen/papers/aws2003.pdf}.
\textit{Formulation source.}
\end{itemize}
\end{document}
